\section{功率信号}\label{sec:Power_Signal}

功率信号是满足$0<P_f<\infty$的信号，其中
\begin{equation}
    P_f=\lim_{T\to\infty}\frac{1}{T}\int_{-T/2}^{T/2}|f(t)|^2\,\dd t.
\end{equation}

为了便于线性运算，下面研究具体的信号空间时也包含零信号，但功率信号并不因此构成一个线性空间，因为两个功率信号之和的功率未必有定义。有限能量信号的平均功率为零；反过来，零平均功率并不保证能量有限。例如，$f(x)=0,\ x<1;\ f(x)=x^{-0.5},\ x\ge 1$，则$\frac{1}{T}\int_{-T/2}^{T/2}|f(x)|^2\,\dd x$与$\frac{\ln T}{T}$同阶。

周期信号是功率信号的重要例子。若函数$f$的周期为$T_0$，且一个周期内的能量
\begin{equation}
    E_{T_0}=\int_0^{T_0}|f(t)|^2\,\dd t
\end{equation}
满足$0<E_{T_0}<\infty$，则
\begin{equation}
    P_f=\frac{E_{T_0}}{T_0}.
\end{equation}

在周期函数的傅里叶分析中，复指数与正弦、余弦是基本成分。有限个这样的成分叠加，既包括有限谐波的周期函数，也包括频率成分之间没有公周期的函数，例如
\begin{equation}
    \sin(2\pi t)+\sin(2\sqrt{2}\pi t).
\end{equation}
这使\textbf{有限频率叠加成为研究功率信号的一类有代表性的对象}。本节以有限个复指数函数的线性组合为主要研究对象，将其组成的空间记为$\mathcal{P}$。我们从时间平均出发，建立它们的互功率、相关函数及功率谱，并说明有限时间截短所得的能量谱如何趋于功率谱。

\subsection{\texorpdfstring{$\mathcal{P}$空间}{P空间}}\label{sub:power_signal}

\begin{definition}
    扩展\cref{def:trinomial}的定义，将任意频率的复指数函数的有限线性组合构成的集合记为
    \begin{equation}
        \mathcal{P}=\braces{
            f(t)=\sum_{k=1}^{N}a_k\e^{2\pi\im\xi_k t}:
            N\in\mathbb{N}^*,\ a_k\in\mathbb{C},\ \xi_k\in\mathbb{R}
        }.
    \end{equation}
\end{definition}

频率$\xi_k$不要求为某个固定基频$\xi_0$的整数倍，否则$f$以$1/\xi_0$为周期。

\begin{proposition}[平均正交性]
    对任意$\xi,\eta\in\mathbb{R}$，有
    \begin{equation}
        \lim_{T\to\infty}\frac{1}{T}\int_{-T/2}^{T/2}
        \e^{2\pi\im\xi t}
        \overline{\e^{2\pi\im\eta t}}\,\dd t
        =\begin{dcases}
            1, & \xi=\eta,   \\
            0, & \xi\ne\eta.
        \end{dcases}
    \end{equation}
\end{proposition}
\begin{proof}
    若$\xi=\eta$，被积函数恒为1。若$\xi\ne\eta$，直接积分得
    \begin{equation}
        \frac{1}{T}\int_{-T/2}^{T/2}\e^{2\pi\im(\xi-\eta)t}\,\dd t
        =\frac{\sin\bigl(\pi(\xi-\eta)T\bigr)}{\pi(\xi-\eta)T}\longrightarrow 0.
    \end{equation}
\end{proof}

下面验证$\mathcal{P}$是功率信号中的线性空间：$\forall f,g\in\mathcal{P}$，取两者的频率的并集$\Lambda$，则
\begin{align}
    \lim_{T\to\infty}\frac{1}{T}\int_{-T/2}^{T/2}f(t)\overline{g(t)}\,\dd t
     & =\sum_{\xi,\eta\in\Lambda}a_\xi\overline{b_\eta}
    \lim_{T\to\infty}\frac{1}{T}\int_{-T/2}^{T/2}
    \e^{2\pi\im(\xi-\eta)t}\,\dd t=\sum_{\xi\in\Lambda}a_\xi\overline{b_\xi}.
\end{align}
特别地，
\begin{equation}
    P_f=\sum_{\xi\in\Lambda}|a_\xi|^2.
\end{equation}

因此，$\mathcal{P}$中的每个非零信号都有正而有限的平均功率，是功率信号。在合并同频率项后，不同频率成分的功率可以相加。

将能量与功率的概念推广到两个信号之间：
\begin{definition}[互能量与互功率]\label{def:cross_energy_power}
    如果$f,g\in L^2(\mathbb{R})$，则它们的\textbf{互能量}定义为：
    \begin{equation}
        E_{fg}=(f,g)=\int_{\mathbb{R}}f\cdot \overline{g}\,\dd m,
    \end{equation}
    如果$f,g\in \mathcal{P}$，则它们的\textbf{互功率}定义为：
    \begin{equation}
        P_{fg}=(f,g)_{\pwr}=\lim_{T\to\infty}\frac{1}{T}\int_{-T/2}^{T/2}f\cdot\overline{g}\,\dd m.
    \end{equation}
\end{definition}

互能量积分的存在性由$L^2$中的柯西—施瓦茨不等式保证；互功率极限的存在性则已由平均正交性证明。

在$\mathcal{P}$上定义\textbf{内积}：
\begin{equation}
    (f,g)_{\pwr}=P_{fg}
    =\sum_{\xi\in\Lambda}a_\xi\overline{b_\xi}.
\end{equation}
这是系数向量的内积：它对第一变元线性、对第二变元共轭线性，满足共轭对称性，并且
\begin{equation}
    (f,f)_{\pwr}=\sum_{\xi\in\Lambda}|a_\xi|^2\ge 0,
\end{equation}
等号成立当且仅当$f=0$。所以它是$\mathcal{P}$上的内积，诱导范数
\begin{equation}
    \|f\|_{\pwr}=\sqrt{P_f}
    =\paren{\sum_{\xi\in\Lambda}|a_\xi|^2}^{1/2}.
\end{equation}

柯西—施瓦茨不等式给出
\begin{equation}
    |P_{fg}|\le\sqrt{P_fP_g}.
\end{equation}

\subsection{功率信号的相关}\label{sub:power_corelation}

\begin{definition}[功率信号的相关]\label{def:power_correlation}
    功率信号$f,g\in \mathcal{P}$的互相关函数为
    \begin{equation}
        R_{fg}(x)=\lim_{T\to\infty}\frac{1}{T}\int_{-T/2}^{T/2}f(x+y)\overline{g(y)}\,\dd y.
    \end{equation}
    当$g=f$时，称为$f$的自相关函数，记为$R_f$。
\end{definition}

下面验证这些极限存在。将两个信号在共同的有限频率集$\Lambda$上写为
\begin{equation}
    f(t)=\sum_{\xi\in\Lambda}a_\xi \e^{2\pi\im\xi t},\qquad g(t)=\sum_{\xi\in\Lambda}b_\xi \e^{2\pi\im\xi t}.
\end{equation}

对每个固定的$\tau$，平移后的信号仍属于$\mathcal{P}$。应用上一小节的平均正交性，得到
\begin{equation}
    R_{fg}(\tau)=\sum_{\xi\in\Lambda}a_\xi\overline{b_\xi}\e^{2\pi\im\xi\tau},\quad R_f(\tau)=\sum_{\xi\in\Lambda}|a_\xi|^2\e^{2\pi\im\xi\tau}.
\end{equation}
由此不仅证明了极限存在，也给出了相关函数的计算方法。不同频率之间的乘积在时间平均后消失；共同频率的乘积则保留下来，并随时延改变相位。

\begin{proposition}[功率信号的相关函数的性质]
    \begin{enumerate}
        \item $R_{fg}(0)=P_{fg},\quad R_f(0)=P_f$；
        \item $R_{fg}(\tau)=\overline{R_{gf}(-\tau)}$，特别地，自相关满足$R_f(-\tau)=\overline{R_f(\tau)}$；若$f$为实函数，则$R_f$也是实值的，因而为偶函数；
        \item
              \begin{equation}
                  |R_{fg}(\tau)|
                  \le\sum_{\xi\in\Lambda}|a_\xi b_\xi|
                  \le
                  \paren{\sum_{\xi\in\Lambda}|a_\xi|^2}^{1/2}
                  \paren{\sum_{\xi\in\Lambda}|b_\xi|^2}^{1/2}
                  =\sqrt{P_fP_g}.
              \end{equation}
              特别地，$|R_f(\tau)|\le R_f(0)=P_f$。
    \end{enumerate}
\end{proposition}

\subsection{能量谱与功率谱}\label{sub:power_spectrum}

简便起见，本小节中，我们将使用频率形式的傅里叶变换。

对于能量信号$f\in L^2(\mathbb{R})$，如果它还是绝对可积的，即$f\in L^1(\mathbb{R})$，则它的傅里叶变换$F(\xi)$存在，并且满足帕塞瓦尔定理（\cref{thm:plancherel}）：
\begin{equation}
    \int_{-\infty}^{\infty}|f(t)|^2\,\dd t=\int_{-\infty}^{\infty}|F(\xi)|^2\,\dd\xi,
\end{equation}
\begin{equation}
    E_{fg}=\int_{\mathbb{R}}f(t)\overline{g(t)}\,\dd t=\int_{\mathbb{R}}F(\xi)\overline{G(\xi)}\,\dd\xi.
\end{equation}

因此，$f\in L^1(\mathbb{R})\cap L^2(\mathbb{R})$的\textbf{能量可以看作是分布在各个频率上的}。类比概率密度函数，我们定义：

\begin{definition}[能量谱密度]
    设能量信号$f,g\in L^1(\mathbb{R})\cap L^2(\mathbb{R})$，则$f$的\textbf{能量谱密度}(energy spectral density, ESD)为
    \begin{equation}
        S^{\mathrm{en}}_f(\xi)=|F(\xi)|^2,
    \end{equation}
    $f$与$g$的\textbf{互能量谱密度}(cross energy spectrum density, CESD)为
    \begin{equation}
        S^{\mathrm{en}}_{fg}(\xi)=\mathcal{F}f(\xi)\overline{\mathcal{F}g(\xi)}.
    \end{equation}
\end{definition}
前者表示能量在频率上的分布，其积分是总能量；后者的积分是互能量，通常为复数。

对于$f\in \mathcal{P}$，为了使用帕塞瓦尔定理，可以考虑功率信号的截短。本小节中，总是用记号$f_T$来表示时域功率信号$f$的截短：
\begin{equation}
    f_T(t)=f(t)\chi_{[-T/2,T/2]}.
\end{equation}
由于$f$有界，$f_T(t)\in L^1(\mathbb{R})\cap L^2(\mathbb{R})$，因此可以定义它的傅里叶变换，记为$F_T(\xi)$。对$f_T(t)$使用帕塞瓦尔定理，有：
\begin{equation}
    E_{f_T}=\int_{\mathbb{R}}|f_T(t)|^2\,\dd t=\int_{\mathbb{R}}|F_T(\xi)|^2\,\dd\xi.
\end{equation}

这样，在$T\to\infty$的极限存在时，$E_{f_T}$的时间平均值就是$f$的功率：
\begin{equation}
    P_f=\lim_{T\to\infty}\frac{1}{T}\int_{-T/2}^{T/2}|f(t)|^2\,\dd t=\lim_{T\to\infty}\frac{E_{f_T}}{T}=\lim_{T\to\infty}\frac{1}{T}\int_{\mathbb{R}}|F_T(\xi)|^2\,\dd\xi.
\end{equation}
如果能交换极限与积分的顺序，则
\begin{equation}
    P_f=\int_{\mathbb{R}}\lim_{T\to\infty}\frac{|F_T(\xi)|^2}{T}\,\dd\xi.
\end{equation}

然而，这样定义的功率谱密度不仅计算流程复杂，而且极限$\displaystyle \lim_{T\to\infty}\frac{1}{T}|F_T(\xi)|^2$未必存在，积分与极限的换序也未必合法。

前面已经说明，$\mathcal{P}$中函数的相关仍属于$\mathcal{P}$，其傅里叶变换（在分布的意义下）是存在的，因此，我们不妨直接扩展帕塞瓦尔定理的适用范围，定义：

\begin{definition}[功率谱密度]
    设功率信号$f,g\in\mathcal{P}$，则$f$的\textbf{功率谱密度}(power spectral density, PSD)为
    \begin{equation}
        S_f=\mathcal{F}R_f,
    \end{equation}
    $f$与$g$的\textbf{互功率谱密度}(cross power spectrum density, CPSD)为
    \begin{equation}
        S_{fg}=\mathcal{F}R_{fg}.
    \end{equation}
\end{definition}

由相关函数的有限频率表达式，可以直接得到功率谱的计算公式。

\begin{theorem}[有限频率信号的功率谱]\label{thm:PSD}
    设$f,g\in\mathcal{P}$，在共同的有限频率集$\Lambda$上写成
    \begin{equation}
        f(t)=\sum_{\xi\in\Lambda}a_\xi\e^{2\pi\im\xi t},\qquad g(t)=\sum_{\xi\in\Lambda}b_\xi\e^{2\pi\im\xi t},
    \end{equation}
    其中，相同频率的项已经合并，未出现的频率对应的系数取零。则
    \begin{equation}
        S_f=\sum_{\xi\in\Lambda}|a_\xi|^2\delta_\xi,\qquad S_{fg}=\sum_{\xi\in\Lambda}a_\xi\overline{b_\xi}\delta_\xi.
    \end{equation}
\end{theorem}
\begin{proof}
    由上一小节，
    \begin{equation}
        R_{fg}(\tau)=\sum_{\xi\in\Lambda}a_\xi\overline{b_\xi}\e^{2\pi\im\xi\tau}.
    \end{equation}
    对这个有限和逐项作傅里叶变换，利用$\mathcal{F}\sbr{\e^{2\pi\im\xi\tau}}=\delta_\xi$，便得
    \begin{equation}
        S_{fg}=\mathcal{F}R_{fg}=\sum_{\xi\in\Lambda}a_\xi\overline{b_\xi}\delta_\xi.
    \end{equation}
    令$g=f$，即得自功率谱公式。
\end{proof}

因此，$\mathcal{P}$中信号的功率谱由有限条冲激谱线组成。频率$\xi$处谱线的强度$|a_\xi|^2$，就是该频率分量的平均功率；各条谱线的强度之和为
\begin{equation}
    \sum_{\xi\in\Lambda}|a_\xi|^2=P_f.
\end{equation}
这说明功率谱确实描述了功率在不同频率之间的分配。类似地，互功率谱中各条谱线的系数之和为$P_{fg}$，但这些系数通常是复数。

沿用前面冲激积分的记法，也可写成
\begin{equation}
    \int_{\mathbb{R}}S_f(\xi)\,\dd\xi=P_f,\qquad \int_{\mathbb{R}}S_{fg}(\xi)\,\dd\xi=P_{fg}.
\end{equation}
这里的积分表示有限个冲激项的系数之和，并不是把$\delta$当作普通可积函数。若用测试函数严格表述，可取一个在全部谱线所在区间上恒为$1$、在区间外光滑过渡到$0$的紧支集函数$\rho$；上面的两个积分分别表示$\langle S_f,\rho\rangle$和$\langle S_{fg},\rho\rangle$，其值与这样的$\rho$的具体选取无关。

注意，虽然$f$本身的傅里叶变换为
\begin{equation}
    \mathcal{F}f=\sum_{\xi\in\Lambda}a_\xi\delta_\xi,
\end{equation}
但不能将功率谱理解为分布$\mathcal{F}f$的模平方。\textbf{将谱线系数取模平方的规则来自平均正交性和相关函数的计算}，不涉及没有定义的$\delta^2$。

\subsubsection{截短功率谱的分布极限}

下面证明，刚才定义的功率谱也是有限时间截短所得能量谱除以窗口长度后的分布极限，从而将相关函数的定义与实际的有限时间观察联系起来。

\begin{theorem}[功率信号的维纳-辛钦定理]\label{thm:power_winner_khinchin}
    设$f,g\in\mathcal{P}$，记
    \begin{equation}
        I_T=[-T/2,T/2],\qquad f_T=f\chi_{I_T},\quad g_T=g\chi_{I_T},
    \end{equation}
    以及$F_T=\mathcal{F}f_T$、$G_T=\mathcal{F}g_T$。则在缓增分布的意义下，
    \begin{equation}
        \lim_{T\to\infty}\frac{F_T\overline{G_T}}{T}=S_{fg}=\mathcal{F}R_{fg}.
    \end{equation}
    特别地，
    \begin{equation}
        \lim_{T\to\infty}\frac{|F_T|^2}{T}=S_f.
    \end{equation}
\end{theorem}
\begin{proof}
    先考察截短信号的互相关函数除以窗口长度所得的函数
    \begin{align}
        q_T(\tau)
         & =\frac{1}{T}\int_{\mathbb{R}}f_T(t+\tau)\overline{g_T(t)}\,\dd t\nonumber    \\
         & =\frac{1}{T}\int_{I_T\cap(I_T-\tau)}f(t+\tau)\overline{g(t)}\,\dd t.
    \end{align}
    这里必须同时保留两个截短窗口。它与相关函数定义中出现的有限平均
    \begin{equation}
        c_T(\tau)=\frac{1}{T}\int_{I_T}f(t+\tau)\overline{g(t)}\,\dd t
    \end{equation}
    一般并不相等。

    $f,g\in\mathcal{P}$有界。固定$\tau$，当$T>|\tau|$时，两者积分区间相差的长度为$|\tau|$，所以
    \begin{equation}
        |q_T(\tau)-c_T(\tau)|\le\frac{|\tau|}{T}\|f\|_\infty\|g\|_\infty\longrightarrow 0.
    \end{equation}
    前面已经证明$c_T(\tau)\to R_{fg}(\tau)$，因此对每个固定的$\tau$，都有$q_T(\tau)\to R_{fg}(\tau)$。

    此外，对所有$T>0$和$\tau\in\mathbb{R}$，
    \begin{equation}
        |q_T(\tau)|\le\|f\|_\infty\|g\|_\infty.
    \end{equation}
    任取$\psi\in\mathcal{S}(\mathbb{R})$，函数$\|f\|_\infty\|g\|_\infty|\psi|$可积。由控制收敛定理（\cref{thm:DCT}），
    \begin{equation}
        \lim_{T\to\infty}\int_{\mathbb{R}}q_T(\tau)\psi(\tau)\,\dd\tau=\int_{\mathbb{R}}R_{fg}(\tau)\psi(\tau)\,\dd\tau.
    \end{equation}
    这就证明了$q_T\to R_{fg}$在缓增分布的意义下成立。

    另一方面，$f_T,g_T\in L^1(\mathbb{R})\cap L^2(\mathbb{R})$，由能量信号的维纳-辛钦定理（\cref{thm:winner-khinchin}），
    \begin{equation}
        \mathcal{F}q_T=\frac{F_T\overline{G_T}}{T}.
    \end{equation}
    再由分布的傅里叶变换的连续性（\cref{thm:distribution_fourier_continuous}），得到
    \begin{equation}
        \frac{F_T\overline{G_T}}{T}=\mathcal{F}q_T\longrightarrow\mathcal{F}R_{fg}=S_{fg}.
    \end{equation}
\end{proof}

\subsubsection{功率谱的计算例子}

\begin{example}[单个余弦函数]
    设$f(t)=A\cos(2\pi\xi_0t+\theta)$，其中$A\in\mathbb{R}$、$\xi_0>0$。由
    \begin{equation}
        f(t)=\frac{A}{2}\e^{\im\theta}\e^{2\pi\im\xi_0t}+\frac{A}{2}\e^{-\im\theta}\e^{-2\pi\im\xi_0t},
    \end{equation}
    得到
    \begin{equation}
        P_f=\frac{A^2}{2},\qquad R_f(\tau)=\frac{A^2}{2}\cos(2\pi\xi_0\tau),
    \end{equation}
    \begin{equation}
        S_f=\frac{A^2}{4}\bigl(\delta_{\xi_0}+\delta_{-\xi_0}\bigr).
    \end{equation}
    初相位不影响自功率谱；两个频率处的谱线强度之和为$A^2/2$，正好等于平均功率。
\end{example}

为了具体说明截短窗口的影响，仍取$f(t)=\cos(2\pi t)$，记截短信号的自相关为
\begin{equation}
    R_T(\tau)=\int_{\mathbb{R}}f_T(t+\tau)f_T(t)\,\dd t.
\end{equation}
当$|\tau|>T$时，两个窗口不相交，故$R_T(\tau)=0$。当$|\tau|\le T$时，交集的长度为$T-|\tau|$；利用积化和差公式积分，得到
\begin{equation}
    R_T(\tau)=\frac{\sin\bigl(2\pi(T-|\tau|)\bigr)}{4\pi}+\frac{T-|\tau|}{2}\cos(2\pi\tau).
\end{equation}
于是对每个固定的$\tau$，
\begin{equation}
    \frac{R_T(\tau)}{T}\longrightarrow\frac{1}{2}\cos(2\pi\tau)=R_f(\tau),
\end{equation}
与上述定理的证明一致。

\begin{example}
    设$f(t)=\sin(2\pi t)+\sin(2\sqrt{2}\pi t)$。由于两个正频率之比为无理数，$f$不是周期函数，但仍属于$\mathcal{P}$。其傅里叶变换为
    \begin{equation}
        \mathcal{F}f=\frac{1}{2\im}(\delta_1-\delta_{-1})+\frac{1}{2\im}(\delta_{\sqrt{2}}-\delta_{-\sqrt{2}}).
    \end{equation}
    因而
    \begin{equation}
        S_f=\frac{1}{4}\bigl(\delta_1+\delta_{-1}+\delta_{\sqrt{2}}+\delta_{-\sqrt{2}}\bigr),\qquad P_f=1.
    \end{equation}
\end{example}

\subsection{为什么不研究完整的功率信号空间}\label{sub:power_signal_scope}

把全部功率信号作为统一的研究对象，则该空间缺少线性空间的结构，因为它对线性组合不封闭；此外，单个信号的平均功率存在，并不足以保证互功率、相关函数等运算也有定义。

\begin{example}
    令$A_n=2^{2^n}$（$n=0,1,\ldots$），取$f(t)=1$，并定义实值偶函数
    \begin{equation}
        g(t)=
        \begin{dcases}
            1,      & |t|<A_0,                              \\
            (-1)^n, & A_n\le|t|<A_{n+1},\quad n=0,1,\ldots.
        \end{dcases}
    \end{equation}
    由于$|f(t)|=|g(t)|=1$，两个信号在任意窗口上的平均功率都等于$1$。然而，$g$在越来越长的区间上交替取$1$和$-1$。取$T_n=2A_{n+1}$，令
    \begin{equation}
        M_n=\frac{1}{T_n}\int_{-T_n/2}^{T_n/2}f(t)\overline{g(t)}\,\dd t=\frac{1}{A_{n+1}}\int_0^{A_{n+1}}g(t)\,\dd t.
    \end{equation}
    最后一段区间的符号为$(-1)^n$，前面所有区间的总长度只有$A_n$，故
    \begin{equation}
        |M_n-(-1)^n|\le\frac{2A_n}{A_{n+1}}\to 0,\quad n\to\infty.
    \end{equation}
    因此互功率沿两条子列分别趋于$1$和$-1$，极限不存在。可见柯西—施瓦茨不等式只能说明这些平均值的模不超过$1$，不能保证它们收敛。
    \begin{equation}
        \frac{1}{T}\int_{-T/2}^{T/2}|f(t)+g(t)|^2\,\dd t=\frac{1}{T}\int_{-T/2}^{T/2}(f^2(t)+2f(t)g(t)+g^2(t))\,\dd t=2+\frac{2}{T}\int_{-T/2}^{T/2}g(t)\,\dd t,
    \end{equation}
    所以$f+g$的自功率平均也没有极限。
\end{example}

相比之下，$\mathcal{P}$中的有限求和与平均正交性足以保证本节各个极限存在，并且已经涵盖周期信号的有限谐波叠加和没有公周期的多频率叠加。对一个周期内平方可积的周期信号，其有限傅里叶和还能在平均平方意义下逼近原信号。这些对象足以呈现功率谱分析的基本方法。

若要进一步允许平均平方意义下的极限，可以按平均功率范数对$\mathcal{P}$作完备化，得到Besicovitch平方平均几乎周期空间，常记为$B^2$；为了定义其中的范数和内积，需要将平均平方距离为零的函数视为等价类。本节不继续展开这些空间的构造。
